\section{作品难点与创新}

本作品面向家庭储能与分布式光伏离网供电场景，提出一种基于 GD32 系列 MCU 的光储逆变一体化能源电力系统。系统由光伏输入模块、储能电池及 BMS 模块、前级三端口功率变换器、后级单相逆变模块和无线监控模块构成，能够完成光伏能量接入、电池安全管理、交流输出供电和远程监测控制。

系统采用 GD32G553、GD32C113 和 GD32VW553 三颗 MCU 分工协作。其中，GD32G553 作为主功率控制器，负责光伏输入功率变换、储能充放电控制、直流母线调节和逆变输出控制；GD32C113 作为电池管理控制器，负责电芯采样、均衡、SOC 估算和故障保护；GD32VW553 作为联网控制器，负责无线通信、参数配置、运行日志和告警信息管理。该架构既保证底层功率环路的实时性，又提升了系统安全性和用户侧交互能力。

\subsection{应用场景与作品定位}

户用光储系统通常需要同时满足光伏高效利用、储能安全运行、离网交流供电和远程运维管理等需求。传统分立式方案往往由光伏控制器、充电器、逆变器和监控模块独立构成，系统体积较大、通信链路复杂、能量调度不够统一。本作品将光伏输入、储能电池和逆变输出统一纳入数字电源控制框架，重点解决小功率家庭备用供电场景下的集成化、智能化和国产化控制实现问题。

本系统的目标不是单一变换器验证，而是完成从电源硬件、控制算法、电池保护到无线监控的整机方案。论文后续章节将围绕系统架构、主功率电路、控制策略、保护机制和测试验证展开。

\subsection{难点分析}

实现一个光伏、储能、逆变一体化的微型数字电源系统，主要面临以下难题：

\begin{enumerate}[label=\arabic*、, leftmargin=*, itemsep=0.8ex, parsep=0pt, topsep=0.5ex]
  \item 小型化趋势与直流母线高压需求之间存在矛盾。为实现兼具体积和功率密度的三端口功率变换器，非隔离拓扑结构成为优选方案。然而，将低压光伏与电池端口提升至逆变所需直流母线，会使变换器在高电压增益条件下运行，对开关器件、电感设计和控制稳定性提出更高要求。
  \item 多端口强耦合与协调控制难。光伏、储能和直流母线三者之间存在能量耦合，不同工况下需要同时满足光伏高效利用、母线稳定、电池安全和负载供电，控制目标相互约束，容易出现能量冲突、母线波动或模式误判。
  \item 多模式平滑切换难度大。系统需在光伏优先、光伏充电、电池放电、混合供电、欠压保护和过载保护等模式间切换，不同模式的控制目标和限值条件差异较大，若切换逻辑处理不当，可能引发电压跌落、电流冲击和逆变输出畸变。
  \item GD32 数字控制实时性要求高。三端口变换和单相逆变需要多路电压、电流和功率同步采样，同时完成 PWM 更新、多环路闭环计算和故障中断处理，对 GD32G553 的外设配置、中断优先级和算法执行效率提出较高要求。
  \item BMS 与主功率系统协同难。电池 SOC 估算、均衡和保护需要与三端口充放电策略深度联动，既要保证电池不过充、不过放和不过温，又不能牺牲母线动态响应和负载供电连续性。
  \item 逆变输出与直流母线动态耦合明显。交流负载突变会直接扰动直流母线，前级三端口调节与后级逆变稳压之间必须协调设计，才能兼顾母线稳定、输出电压质量、动态响应和带载能力。
  \item 多 MCU 协同与无线监控实时性需要平衡。GD32G553、GD32C113 和 GD32VW553 分工协作，需要设计可靠通信协议和故障联动逻辑；无线监控数据上传和参数下发不能干扰底层功率控制实时性。
\end{enumerate}

\subsection{创新点分析}

针对以上难点，本团队从拓扑集成、控制架构、能量调度和系统安全四个层面进行设计，形成以下创新点：

\begin{enumerate}[label=\arabic*、, leftmargin=*, itemsep=0.8ex, parsep=0pt, topsep=0.5ex]
  \item 非隔离三端口光储变换架构。系统将光伏端口、储能端口和直流母线端口统一到前级功率变换器中，相比传统分立式光伏控制器加电池充电器方案，减少能量变换级数和控制接口，有利于提升整机集成度和响应速度。
  \item 多 MCU 协同控制架构。构建 GD32G553、GD32C113、GD32VW553 三芯片协同架构，分别承担主功率控制、电池安全管理和无线监控任务，实现实时控制、安全保护和通信交互解耦，充分发挥 GD32 系列 MCU 在数字电源控制、BMS 管理和无线互联方面的综合能力。
  \item 光储协同多模态能量调度策略。提出以光伏优先利用、母线稳定和电池安全为约束的能量调度逻辑，根据光伏功率、负载功率、电池 SOC 和故障状态自动选择工作模式，实现储能自适应充放电和离网稳压运行。
  \item 面向工程应用的分层保护机制。系统在硬件比较、MCU 快速中断、BMS 状态机和无线告警层面设置多级保护，对过压、欠压、过流、过温、通信异常和电池异常进行分层处理，提高整机运行可靠性。
  \item 可扩展远程监控与参数配置。基于 GD32VW553 实现无线联网，将运行数据、故障日志和关键参数上传至上位机或移动端，为后续远程运维、数据分析和算法在线优化提供接口。
\end{enumerate}
